% ACHTUNG: Für das Erstellen des Literaturverzeichnisses wird das modernere Paket biblatex
%			 in Kombination mit biber verwendet -- nicht mehr das ältere BibTex!
% 			 Bitte stellen Sie ggf. Ihre TeX-Umgebung
% 			 entsprechend ein (z.B. TeXStudio: Einstellungen --> Erzeugen --> Standard Bibliographieprogramm: biber)
%

\documentclass[
	12pt,
	BCOR=5mm,
	DIV=12,
	headinclude=on,
	footinclude=off,
	parskip=half,
	bibliography=totoc,
	listof=entryprefix,
	toc=listof,
	pointlessnumbers,
	plainfootsepline]{scrreprt}

%	Konfigurationsdatei einziehen
% !TEX root =  master.tex
% 		HYPERREF
%
\usepackage[
	hidelinks=true % keine roten Markierungen bei Links
]{hyperref}

% Zwei eigene Befehle zum Setzen von Autor und Titel. Ausserdem werden die PDF-Informationen richtig gesetzt.
\newcommand{\TitelDerArbeit}[1]{\def\DerTitelDerArbeit{#1}\hypersetup{pdftitle={#1}}}
\newcommand{\AutorDerArbeit}[1]{\def\DerAutorDerArbeit{#1}\hypersetup{pdfauthor={#1}}}
\newcommand{\Firma}[1]{\def\DerNameDerFirma{#1}}
\newcommand{\Kurs}[1]{\def\DieKursbezeichnung{#1}}


% Correct superscripts 
\usepackage{fnpct}


%		FONT AND INPUT ENCODING
%
\usepackage[T1]{fontenc}
\usepackage[utf8]{inputenc}

%		CALCULATIONS
%
\usepackage{calc} % Used for extra space below footsepline

%		LANGUAGE SETTINGS
%
\usepackage[ngerman]{babel} 	% German language
\usepackage[german=quotes]{csquotes} 	% correct quotes using \enquote{}

%\usepackage[english]{babel}   % For english language
%\usepackage{csquotes} 	% Richtiges Setzen der Anführungszeichen mit \enquote{}


%		BIBLIOGRAPHY SETTINGS
%

% Uncomment the next three lines for author-year-style with footnotes (Chicago)
\usepackage[backend=biber, autocite=footnote, style=authoryear, dashed=false]{biblatex} 	%Use Author-Year-Cites with footnotes
\AdaptNoteOpt\footcite\multfootcite   %will add  separators if footcite is called multiple consecutive times 
\AdaptNoteOpt\autocite\multautocite % will add  separators if autocite is called multiple consecutive times

% Uncomment the next line for IEEE-style 
% \usepackage[backend=biber, autocite=inline, style=ieee]{biblatex} 	% Use IEEE-Style (e.g. [1])

% Uncomment the next line for alphabetic style 
% \usepackage[backend=biber, autocite=inline, style=alphabetic]{biblatex} 	% Use alphabetic style (e.g. [TGK12])

% Uncomment the next two lines vor Harvard-Style 
%\usepackage[backend=biber, style=apa]{biblatex} 	
%\DeclareLanguageMapping{german}{german-apa}


\DefineBibliographyStrings{ngerman}{  %Change u.a. to et al. (german only!)
	andothers = {{et\,al\adddot}},
}

%%% Uncomment the following lines to support hard URL breaks in bibliography 
%\apptocmd{\UrlBreaks}{\do\f\do\m}{}{}
%\setcounter{biburllcpenalty}{9000}% Kleinbuchstaben
%\setcounter{biburlucpenalty}{9000}% Großbuchstaben


\setlength{\bibparsep}{\parskip}		%add some space between biblatex entries in the bibliography
\addbibresource{.bib}	%Add file bibliography.bib as biblatex resource

%		CORRECT PAGE MARGINS
%			ADDED BY JANIK STRACKE
\usepackage[a4paper,
left=35mm,% Left Margin: 3.5cm
right=25mm,% Right Margin: 2.5cm
top=15mm,% Top Margin to head: 1.5cm
bottom=15mm,% Bottom Margin to foot: 1.5cm
includehead,
includefoot]{geometry}

%		SET LINE SPACING TO ONE HALF
%			ADDED BY JANIK STRACKE
\usepackage[onehalfspacing]{setspace}


%		FOOTNOTES 
%
% Count footnotes over chapters
\usepackage{chngcntr}
\counterwithout{footnote}{chapter}

%	ACRONYMS
%%%
%%% WICHTIG: Installieren Sie das neueste Acronyms-Paket!!!
%%%
\makeatletter
\usepackage{acronym}
% \@ifpackagelater{acronym}{2015/03/20}
%{%
% \renewcommand*{\aclabelfont}[1]{\textbf{\textsf{\acsfont{#1}}}}
%}%
%{%
%}%
\makeatother

%		LISTINGS
\usepackage{listings}	%Format Listings properly
\renewcommand{\lstlistingname}{Quelltext} 
\renewcommand{\lstlistlistingname}{Quelltextverzeichnis}
\lstset{numbers=left,
	numberstyle=\tiny,
	captionpos=b,
	basicstyle=\ttfamily\small}


%		EXTRA PACKAGES

\usepackage{lipsum}    %Blindtext
\usepackage{graphicx} % use various graphics formats
\usepackage[german]{varioref} 	% nicer references \vref
\usepackage{caption}	%better Captions
\usepackage{booktabs} %nicer Tabs
\usepackage{array}
%\newcolumntype{P}[1]{>{\raggedright\arraybackslash}p{#1}}


%		ALGORITHMS
\usepackage{algorithm}
\usepackage{algpseudocode}
\renewcommand{\listalgorithmname}{Algorithmenverzeichnis}
\floatname{algorithm}{Algorithmus}


%		FONT SELECTION: Entweder Latin Modern oder Times / Helvetica
\usepackage{lmodern} %Latin modern font
%\usepackage{mathptmx}  %Helvetica / Times New Roman fonts (2 lines)
%\usepackage[scaled=.92]{helvet} %Helvetica / Times New Roman fonts (2 lines)


%		PAGE HEADER / FOOTER
%	    Warning: There are some redefinitions throughout the master.tex-file!  DON'T CHANGE THESE REDEFINITIONS!
\RequirePackage[automark,headsepline,footsepline]{scrpage2}
\pagestyle{scrheadings}
\renewcommand*{\pnumfont}{\upshape\sffamily}
\renewcommand*{\headfont}{\upshape\sffamily}
\renewcommand*{\footfont}{\upshape\sffamily}
\renewcommand{\chaptermarkformat}{}
\RedeclareSectionCommand[beforeskip=0pt]{chapter}
\clearscrheadfoot

\ifoot[\rule{0pt}{\ht\strutbox+\dp\strutbox}SAP SE / DHBW Mannheim]{\rule{0pt}{\ht\strutbox+\dp\strutbox}SAP SE / DHBW Mannheim}
\ofoot[\rule{0pt}{\ht\strutbox+\dp\strutbox}\pagemark]{\rule{0pt}{\ht\strutbox+\dp\strutbox}\pagemark}

\ohead{\headmark}

\usepackage{pdfpages}

\begin{document}

%% BITTE GEBEN SIE HIER DEN TITEL UND DIE AUTORIN / DEN AUTOR DER ARBEIT AN!
%% DIESE INFORMATIONEN _MÜSSEN_ GESETZT SEIN, UM TITELBLATT, ABSTRACT UND
%% EIGENSTÄNDIGKEITSERKLÄRUNG AUTOMATISCH ANZUPASSEN!
\TitelDerArbeit{Möglichkeiten zum personalisierten Lernen im Studium}
\AutorDerArbeit{Andrea Engel}
\Firma{SAP SE}
\Kurs{WWI 17 SEB}

\begin{titlepage}
\begin{minipage}{\textwidth}
		\vspace{-2cm}
		\noindent \includegraphics[scale=0.14]{img/sap.png} \hfill   \includegraphics{img/logo.jpg}
\end{minipage}
\vspace{1em}
\sffamily
\begin{center}
	\textsf{\large{}Duale Hochschule Baden-W\"urttemberg\\[1.5mm] Mannheim}\\[2em]
	\textsf{\textbf{\Large{}Seminararbeit}}\\[3mm]
	\textsf{\textbf{Xpire - Dokumentation zum Projekthergang}} \\[1.5cm]
	\textsf{\textbf{\Large{}Studiengang Wirtschaftsinformatik}\\[3mm] \textsf{Studienrichtung Software Engineering}}
	
	\vspace{3em}
%	\textsf{\Large{Sperrvermerk}}
\vfill

\begin{minipage}{\textwidth}

\begin{tabbing}
	Wissenschaftlicher Betreuer: \hspace{0.85cm}\=\kill
	Verfasser/in: \> Felix Waage, Jonas Dutzi, Milena Zahn, \\ \> Yvonne Werner, Verena Richarz, Andrea Engel \\[1.5mm]
	Firma: \> SAP SE \\[1.5mm]
	Kurs: \> WWI 17 SEB\\[1.5mm]
	Studiengangsleiter: \> Prof. Dr. Sebastian Ritterbusch  \\[1.5mm]
	Modul: \> Entwicklung mobiler Applikationen \\[1.5mm]
	Dozent/-in: \> Michael Spengler \\[1.5mm]
	Bearbeitungszeitraum: \> 11.05.2020 - 31.07.2020
	
\end{tabbing}
\end{minipage}

\end{center}

\end{titlepage}

\pagenumbering{roman} % Römische Seitennummerierung
\normalfont

%--------------------------------
% Verzeichnisse - nicht benötige Verzeichnisse bitte auskommentieren / löschen.
%--------------------------------

%   Sperrvermerk
%\chapter*{}
Die nachfolgende Arbeit enthält vertrauliche Daten der:

\vspace{0.5cm}

	SAP SE
	
	Dietmar-Hopp-Allee 16
	
	69190 Walldorf
	
\vspace{0.5cm}
Der Inhalt dieser Arbeit darf weder als Ganzes noch in Auszügen Personen außerhalb des Prüfungsprozesses und des Evaluationsverfahrens zugänglich gemacht werden, sofern keine anders lautende Genehmigung der Ausbildungsstätte vorliegt. 
\cleardoublepage


% Ehrenwörtliche Erklärung ewerkl.tex einziehen
% \input{ewerkl.tex}

%	Kurzfassung
% \input{abstract}

%	Inhaltsverzeichnis
\tableofcontents

%	Abbildungsverzeichnis
%\listoffigures

%	Tabellenverzeichnis
% \listoftables

%	Listingsverzeichnis
% \lstlistoflistings

% 	Algorithmenverzeichnis
%\listofalgorithms

% 	Abkürzungsverzeichnis (siehe Datei acronyms.tex!)
% !TEX root =  master.tex

\clearpage
\chapter*{Abkürzungsverzeichnis}	
\addcontentsline{toc}{chapter}{Abkürzungsverzeichnis}


\begin{acronym}
	\acro{DHBW}{Duale Hochschule Baden-Württemberg}

\end{acronym}


%--------------------------------
% Start des Textteils der Arbeit
%--------------------------------
\clearpage
\ihead{\chaptername~\thechapter} % Neue Header-Definition
\pagenumbering{arabic}  % Arabische Seitenzahlen

%	Anleitungs-Datei anleitung.tex einziehen. Auf diese Weise sollten Sie versuchen, für jedes einzelne
% Kapitel eine eigene Datei anzulegen und mittels input-Kommando einzuziehen.

% !TEX root =  master.tex
\chapter{Einleitung}

\section{Grundidee von Xpire}

\section{Was ist eine PWA?}
\input{tex/chapter01/1.1.grundidee}
% !TEX root =  master.tex
\section{Was ist eine PWA?}
Bei der traditionellen Entwicklung mobiler Anwendungen war die Wiederverwendbarkeit von Code zwischen nativen Anwendungen, dem Web und mobilen Plattformen von Natur aus nicht gegeben. Dies ist auf die nicht kompatible Codebasis nativer Anwendungen zurückzuführen, was zu getrennten Projekten und Entwicklerumgebungen führt, wenn Unterstützung für mehrere Plattformen und Betriebssysteme gewünscht wird.\autocite[vgl.][]{Maj}
Progressive Web Apps revolutionieren diesen Ansatz.\\
Als Webseite, die jegliche Eigenschaften einer nativen Applikation aufweist, kann sie auf dem Endgerät über den Webbrowser installiert werden. Diese Eigenschaften beinhalten Installierbarkeit, Offlinefähigkeit oder Push-Benachrichtigungen, selbst wenn die Anwendung geschlossen wurde. Die progressive Web App stellt daher eine Mischung aus responsiver Webseite und nativen App dar.
% Vorgehensweise
Diese plattformübergreifende Web Apps werden auf der Grundlage von HTML5, CSS3 und JavaScript erstellt. Das HTTP-Protokoll zur Kommunikation zwischen Webclient und Webserver stellt dabei die grundlegende Bedingung zur Verwendung einer PWA dar.
Die PWA enthält sogenannte \textit{ServiceWorker}, welche Funktionalitäten wie Offlinefähigkeit, Caching als auch Push-Benachrichtigungen ermöglichen. Bei einem ServiceWorker handelt es sich um einen in JavaScript implementierten Proxy, der zwischen der Anwendung und dem Server geschaltet wird.\autocite[vgl.][]{Hume.2018}
% Vorteile
Die Verwendung einer PWA bietet sowohl für Entwickler als auch für den Anwender viele Vorteile. Durch den Mobile-First-Ansatz wird garantiert, dass die Web Apps problemlos auf mobilen Geräten innerhalb des Browsers funktionieren. Somit ist die User Experience wie in er nativen App und bietet viele bekannte Vorteile während ist die Entwicklung in vielen Fällen um ein Vielfaches schneller, einfacher und kostengünstiger ist als die einer nativen App.


% !TEX root =  master.tex
\section{Vorgehensweise}
Um so effektiv wie möglich zusammen zu arbeiten, haben wir uns eine flexible Art der Aufgabenaufteilung überlegt. Dabei sind wir wie folgt vorgegangen: Zunächst haben wir gemeinsam in einer Brainstorming-Session überlegt, welche Funktionalitäten die Xpire-App beinhalten soll, wie sie aussehen soll und welche Technologien wir für die Umsetzung verwenden möchten. Anschließend hat sich jedes Gruppenmitglied selbstständig mit den Grundlagen einer PWA als auch mit der Softwarebibliothek React vertraut gemacht. In einer zweiten Session haben wir gemeinsam die Code-Basis für das Projekt aufgesetzt und spezifische Verantwortungsbereiche im Projekt definiert, zu denen sich jeder, je nach Stärken und Interessen, selbst zuordnen konnte. Dabei ist zu beachten, dass eine Person nicht zwangsläufig nur im eigenen Verantwortungsbereich agieren muss. Jeder darf entsprechend seiner Stärken den Fokus setzen und darüber hinaus weiteren Aufgaben außerhalb seines Schwerpunktes nachgehen und die anderen Team-Mitglieder bei ihren Aufgaben unterstützen. Wichtig ist uns, dass sich jeder im Team wohl fühlt, dem nachgehen kann, worauf er Lust hat und, ganz wichtig, durch Learning-by-Doing ganz viel Neues erlernen kann.\\
Die Schwerpunkte wurden wie folgt eingeteilt:
\begin{itemize}[noitemsep]
	\item \textbf{Design:} Verena, Andrea
	\item \textbf{App-Entwicklung:} Felix, Milena, Yvonne, Verena
	\item \textbf{Datenbank:} Fabio, Jonas
	\item \textbf{Projektmanagement:} Andrea
\end{itemize}
Freitags, an unseren offiziellen Meeting-Terminen, stellen wir einander vor, welche Aufgaben in der Zwischenzeit umgesetzt wurden, definieren neue Aufgaben und reflektieren die bisherige Vorgehensweise. Unter der Woche vereinbaren wir je nach Bedarf individuell Termine mit dem gesamten Team oder treffen uns in Teilgruppen, um bestimmte Aufgaben umzusetzen.\\
Das gesamte Projekt kennzeichnet sich daher durch eine iterativen Arbeitsweise, sodass genug Platz für neue Ideen und Erweiterungen bleibt. Ziel ist es, zunächst ein Minimal-Viable-Product (MVP) zu entwickeln, um ein funktionierendes Produkt liefern zu können, welche sich in zukünftigen Schritten einfach erweitern lässt.



% !TEX root =  master.tex
\chapter{Analyse}

% !TEX root =  master.tex
\section{Value Proposition Canvas}

% !TEX root =  master.tex
\section{Benutzerprofilanalyse}
Im Rahmen der Benutzerprofilanalyse werden die Eigenschaften der Endanwender untersucht, um die Gebrauchstauglichkeit der Anwendung sicherzustellen. Eine wesentliche
Rolle spielen dabei die persönlichen Merkmale, Fähigkeiten und Kenntnisse der Benutzer
sowie deren Einordnung in ein bestimmtes Nutzungsverhalten.\\
Da die progressive Web App von möglichst vielen Nutzern verwendet werden soll, gibt es kaum Einschränkungen, was die Zielgruppe betrifft. Die Gemeinsamkeit der Nutzer liegt im Besitz eines Kühlschrankes, dem selbstständigem Tätigen von Lebensmittel-Einkäufen sowie im Besitz eines Smartphones oder eines Computers. Daraus lässt sich ableiten, dass die Zielgruppe überwiegend volljährig ist und ein Smartphone oder einen Computer besitzt. Laut Statista besitzen aktuell 66,5 Mio Menschen in Deutschland ein Smartphone.\autocite[vgl.][]{Statista.2020} Ca. 13 Mio Menschen in Deutschland sind Minderjährig, wobei nicht alle Minderjährigen ein Smartphone besitzen. Zieht man trotzdem die Anzahl aller Minderjährigen von der Anzahl der Menschen, die in Deutschland ein Smartphone besitzen, ab, kommt man immer noch auf eine Anzahl von mindestens 49 Mio potentiellen Nutzer. Das sind  mehr als 50\% aller deutschen Staatsbürger.\\
All diese Nutzer können einen völlig unterschiedlichen Erfahrungsschatz bezüglich der Nutzung von neuen Technologien aufweisen. Von der Rentnerin, die nur selten einen Computer verwendet bis hin zum zwanzigjährigen Studierenden, der selbstständig Web Apps programmiert, können die unterschiedlichsten Ausgangsvoraussetzungen gegeben sein. Es is davon auszugehen, dass ein Großteil der Anwendet eine durchschnittliche Technikaffinität aufweist, es aber auch Anwender gibt, die keinerlei oder kaum Erfahrungen im Umgang mit Web Apps aufweisen. Demzufolge sollte die Oberfläche gleichzeitig so konzipiert sein, dass sie auch für Laien intuitiv und verständlich ist. Ein strukturierter und klarer Aufbau könnten diesem Aspekt Rechnung tragen. Schließlich sei noch zu erwähnen, dass die Web App vorerst nur für Deutschland entwickelt wird. Zu einem späteren Zeitpunkt ist die Internationalisierung der Applikation durchaus denkbar, dafür bedarf es dann einer Möglichkeit zur Auswahl unterschiedlicher Sprachen.




% !TEX root =  master.tex
\section{Anforderungsanalyse}
% !TEX root =  master.tex
\section{Monetarisierungsstrategie}
% kostenlos -> viele sollen die App nutzen
% Zusatz: 
%	- Angebote von Supermärkten
%	- Rezeptvorschläge für die Produkte die ablaufen
% Monetarisierung: freiwillige Spenden, Kooperation mit Chefkoch, Supermärkten, restegourmet etc.
% Schalten von werung: darf aber nicht unseriös, aufdringlich oder nervig wirken -> Webrung muss gegen Bezahlung deaktiviert werden können
% nicht alle Möglichkeiten können gleichzeitig umgesetzt werden

% am Ende: keine Werbung -> Finanzierung durch Spenden
% gemeinnütziges Projekt -> Kostenlose Infrastruktur zB bei Strato

% !TEX root =  master.tex
\chapter{Konzeption und Entwurf}

\section{Eingesetzte Technologien}

\section{Konzeptuelles Modell}

\section{Entwicklung von Prototypen}

% !TEX root =  master.tex
\section{Eingesetzte Technologien}


% !TEX root =  master.tex
\newpage
\section{Konzeptuelles Modell}
% Welche Architektur und Warum?
% User-Authentifizierung?
% Welche Persistierung?
% Welche Tests und warum?

In den vorherigen Abschnitten wurden bereits einige Merkmale der Xpire-App beschrieben. Darauf aufbauend sollen im folgenden Abschnitt die Zusammenhänge der verschiedenen Bestandteile näher erläutert werden, um einen detaillierten Überblick über das das konzeptionelle Modell zu liefern.\\
Kern der Xpire-App ist das Frontend, welches mit Hilfe von React als PWA konzipiert ist. Wie den Abbildungen \ref{fig:prot1} und \ref{fig:prot2} entnommen werden kann, besteht die Xpire-App aus 3 Ansichten: \textit{Home-Screen}, \textit{Product-Screen} und dem \textit{Create-Screen}. Diese Ansichten werden in React durch Komponenten realisiert, welche in unterschiedlichen Kontexten wiederverwendet werden können. Der Home-Screen besitzt demzufolge die Komponente AppBar und eine Komponente zur Listendarstellung der Produkte. Die Ansichten \textit{Product-Screen} und \textit{Create-Screen} werden durch dieselbe Komponente realisiert, da zum Erstellen und Anzeigen der Produktinformation ähnliche UI-Bestandteile benötigt werden. Die Darstellung passt sich hierbei automatisch anhand der übergebenen Parameter an. Hierdurch können Code-Duplikate vermieden werden und auch der Raum für Fehler wird reduziert.\\
Ein Weiterer wichtiger Bestandteil der Anwendung ist die IndexedDB, welche zur dauerhaften Persistierung der Daten verwendet wird. Wie bereits in Abschnitt \ref{chapter:datenbank} erwähnt wird hierfür das Modul \textit{Dexie} verwendet. Neben der Datenhaltung in der lokalen Datenbank wird eine aktuelle Kopie der gesamten Datenbank in einem Array im \textit{State} der App-Komponente synchron gehalten, wodurch sowohl die Listen-Komponente als auch die Product-Screen-Komponente darauf zugreifen können.

Der Nutzer soll später seine Produkte mit Hilfe des einheitlichen Barcodes hinzufügen können. Dies soll die Verwendung deutlich beschleunigen und vereinfachen, da so Informationen wie Titel, Gewicht, usw. nicht manuell eingetragen werden müssen. Um die dazu benötigten Informationen zu erhalten, wird die OpenFoodFacts-Api verwendet. Diese führt eine Datenbank mit umfassenden Produktinformationen und kann Anfragen über den Barcode entgegennehmen. Des Weiteren wird diese API verwendet, um Abbildungen der einzelnen Produkte zu erhalten. Dadurch wird dem Nutzer die Verwendung der Anwendung zusätzlich vereinfacht. Darüber hinaus ist in einer späteren Version angedacht, die Produktbilder lokal in der Datenbank zu speichern, um auch bei Offlin-Betrieb die Bilder anzeigen zu können.
% !TEX root =  master.tex
\section{Entwicklung von Prototypen}

% !TEX root =  master.tex
\chapter{Implementierung}


\input{tex/chapter04/4.1.Präsentation}
% !TEX root =  master.tex
\section{Anforderungsabgleich}


% !TEX root =  master.tex
\chapter{Zusammenfassung und Ausblick}
% Was passiert, wenn man ein Produkt löscht? Ist es dann in der Datenbank geschlöscht oder wird nur ein Attribut gesetzt? Welche Auswirkungen hätten diese Möglichkeiten?


%	Literaturverzeichnis
\ihead{} % Neue Header-Definition
\printbibliography
\cleardoublepage

% Der Anhang beginnt hier - jedes Kapitel wird alphabetisch aufgezählt. (Anhang A, B usw.)
\appendix
\ihead{\appendixname~\thechapter} % Neue Header-Definition

% appendix.tex einziehen





\end{document}
